\section{距离含义和可比性}

本章计算目的：说明同一张最终表中不同源项的距离变量含义，并标注哪些模型可按远场距离直接比较，哪些是中心线模型或实验等效场。

\begin{table}[H]
\centering
\caption{辅助表：距离含义与可比性}
\label{tab:distance-meaning}
\small
\begin{tabular}{p{3.5cm}p{4.2cm}p{4.0cm}p{3.0cm}}
\toprule
源项 & 距离含义 & 可比性 & 最终表备注写法\\
\midrule
UEP/ICCP 静态或慢变电场 & $R$：探测器到等效源中心距离 & 可按远场偶极模型给出 100--1000 m 强度 & 远场偶极模型\\
轴频基频 & $R$：探测器到等效源中心距离 & 与静态场同一 $R^{-3}$ 框架下可比较 & 远场偶极调制模型\\
轴频二/三阶谐波 & $R$：探测器到等效源中心距离 & 与轴频基频同一调制框架下可比较 & 远场偶极调制模型\\
尾流/流动诱导电场 & $x$：尾流中心线下游距离 & 使用中心线衰减模型，不等同于径向远场偶极模型 & 尾流中心线模型，非径向远场\\
变频器--水泵系统高频电磁背景 & $r_{\mathrm{lab}}$：实验室 VFD、电机电缆、泵体、水体、接地、DAQ 和传感器链路之间的耦合距离 & 只用频率相关共模/漏电流等效偶极模型描述，不参与真实水下目标远场排序 & 实验 VFD 背景场\\
\bottomrule
\end{tabular}
\end{table}

本报告在目标源项主表中统一列出 100、300、500、1000 m 四个尺度，但其物理含义按表 \ref{tab:distance-meaning} 读取：UEP/轴频使用 $R$，尾流使用 $x$。变频器--水泵系统高频电磁背景单独列为 0--5 m 实验室耦合尺度表，不写入目标源项主表。

本章对客观参考表的作用：为最终表中“距离含义”和“备注”两列提供读取规则，避免把不同距离变量误读为同一种远场传播距离。
